\documentclass{beamer}

\usetheme{metropolis}
\usepackage{plex-sans}

\title{PRIMES in P}
\author{Siddharth \& Amit}
\date{\today}

\begin{document}

\maketitle

\section{Motivation}

\begin{frame}{What is the problem?}
Given a number $n$, determine whether it is prime.

\pause

Naive idea:
\begin{itemize}
    \item Try all divisors up to $\sqrt{n}$
\end{itemize}

\pause

Question:
\begin{itemize}
    \item Is this efficient?
\end{itemize}
\end{frame}

\section{Input Size vs Value}

\begin{frame}{Input Size vs Value}
Computers do not see $n$ directly.

\pause

They see its binary representation:
\[
\text{input size} = \log n
\]

\pause

Trial division takes:
\[
O(\sqrt{n})
\]

\pause

Rewrite in terms of input size:
\[
\sqrt{n} = 2^{\frac{1}{2} \log n}
\]

\pause

This is exponential in input size.
\end{frame}

\begin{frame}{Why this matters}
Polynomial time means:
\[
O((\log n)^k)
\]

\pause

But trial division is:
\[
O(2^{(\log n)/2})
\]

\pause

Conclusion:
\begin{itemize}
    \item Trial division is not in P
\end{itemize}
\end{frame}

\section{Next Step}

\begin{frame}{Where do we go from here?}
We need faster methods.

\pause

Idea:
\begin{itemize}
    \item Use number theory instead of brute force
\end{itemize}

\pause

Leads to:
\begin{itemize}
    \item Fermat test
    \item Miller-Rabin
    \item AKS
\end{itemize}
\end{frame}

\begin{frame}{Takeaway}
\begin{itemize}
    \item Efficiency depends on input size, not value
    \item $\sqrt{n}$ is exponential in $\log n$
    \item We need polynomial-time algorithms in $\log n$
\end{itemize}
\end{frame}

\end{document}